We hereafter survey the existing techniques for detecting anomalies in multi-service applications, \viz symptoms of failures affecting the whole application or some of its services.
We organise the survey by distinguishing three main classes of techniques:
we separate the techniques enacting anomaly detection by directly processing the logs produced by application services (\Cref{sec:detection:service-logs}) from those requiring to instrument applications, \viz to feature distributed tracing (\Cref{sec:detection:distributed-tracing}) or by installing agents to monitor application services (\Cref{sec:detection:monitoring}).
We finally discuss the surveyed anomaly detection techniques, with the support of a recapping table (\Cref{tab:detection}) and by also highlighting some open challenges in anomaly detection (\Cref{sec:detection:discussion}).

\subsection{Log-based Anomaly Detection Techniques}
\label{sec:detection:service-logs}
% The most basic technique for online anomaly detection in multi-service applications is to compare the information logged by the application frontend against the application's SLOs (Service-Level Objectives).
% In case of SLO violations, the whole application is considered to suffer from performance anomalies.
% An alternative solution, which allows detecting anomalies at a finer granularity (\viz detecting anomalous services, rather than anomalies affecting a whole application) is to exploit unsupervised machine learning algorithms to process the logs produced by all application services to detect whether some performance anomaly is occurring in some application services.

% \subsubsection{SLO Check}
% Microscope \cite{06_Lin2018_Microscope,19_Guan2019_MicroscopeDemo} enacts online anomaly detection in multi-service applications by comparing the performances of the application frontend with the SLOs to be enforced by the application.
% In particular, Microscope \cite{06_Lin2018_Microscope,19_Guan2019_MicroscopeDemo} assumes the services forming an application to log their request processing time, \viz the time needed to process each request from an user.
% Microscope \cite{06_Lin2018_Microscope} then compares the request processing time logged by the service acting as the application with that declared in the SLOs of the application. 
% Whenever the performance of the application frontend violates the application SLOs, Microscope \cite{06_Lin2018_Microscope,19_Guan2019_MicroscopeDemo} considers the whole application as affected by a performance anomaly.
Log-based online anomaly detection is typically enacted by exploiting unsupervised machine learning algorithms, which are used to process the logs produced by the services forming an application to detect whether some anomaly is occurring in some application service.

\subsubsection{Unsupervised Learning}
Online anomaly detection in application services can be enacted by exploiting unsupervised learning algorithms to learn a baseline modelling of the logging behaviour of the application services in failure-free runs of an application.
The baseline model is then used online to detect whether the events newly logged by a service diverge from the baseline model, hence indicating that the service suffers from some anomaly.
This technique is adopted by OASIS \cite{38_Nandi2016_AnomalyDetectionControlFlowGraphMining}, Jia \etal \cite{31_Jia2017_AnomalyDiagnosisHybridGraphModel}, and LogSed \etal \cite{32_Jia2017_LogSed}, which we discuss hereafter.

OASIS \cite{38_Nandi2016_AnomalyDetectionControlFlowGraphMining} starts from the application logs generated in failure-free runs of the application, which are used as the training data to mine a control-flow graph modelling the baseline application behaviour under normal conditions, \viz which events should be logged by which service, and in which order.
OASIS \cite{38_Nandi2016_AnomalyDetectionControlFlowGraphMining} mines log templates from the textual content in the training logs, and it clusters the log templates so that all templates originated by the same log statement go in the same cluster.
The nodes in the control flow graph are then determined by picking a representative template from each cluster.
The edges are instead mined based on temporal co-occurrence of log templates, upon which OASIS \cite{38_Nandi2016_AnomalyDetectionControlFlowGraphMining} determines the immediate successors of each template, together with branching probabilities and expected time lags between each template and its successors.
Online anomaly detection is then enacted by mapping newly generated logs to their corresponding templates and by checking whether the arriving log templates adhere to the control-flow graph.
If none of the templates that should follow a template is logged in the corresponding time lag, or if the actual rate of templates following a given template significantly deviates from the expected branching probabilities, OASIS \cite{38_Nandi2016_AnomalyDetectionControlFlowGraphMining} considers the application as affected by a functional anomaly.

Jia \etal \cite{31_Jia2017_AnomalyDiagnosisHybridGraphModel} and LogSed \etal \cite{32_Jia2017_LogSed} provide two other techniques for the detection of functional or performance anomalies in multi-service applications based on the analysis of their logs.
% They process the application logs collected in normal, failure-free runs of a multi-service application to mine a topology graph modelling the dependencies among the application services and, for each service, a time-weighted control flow graph modelling its internal flow.
% By applying the PC-algorithm \cite{Spirtes2000_CausationPredictionSearchPCAlgoritm} to the services' logs, Jia \etal \cite{31_Jia2017_AnomalyDiagnosisHybridGraphModel} and LogSed \etal \cite{32_Jia2017_LogSed} learn a topology graph modelling causal relationships among services, \viz whether some events happening in a service can cause events to happen in other services.
% They then follow a similar technique to that in OASIS \cite{38_Nandi2016_AnomalyDetectionControlFlowGraphMining} to learn the time-weighted control flow graph for each service:
They process the logs collected in normal, failure-free runs of an application to mine, for each service, a time-weighted control flow graph modelling its internal flow.
In particular, Jia \etal \cite{31_Jia2017_AnomalyDiagnosisHybridGraphModel} and LogSed \etal \cite{32_Jia2017_LogSed} learn the time-weighted control flow graph for each service, with a similar method to that in OASIS \cite{38_Nandi2016_AnomalyDetectionControlFlowGraphMining}:
template mining is used to transform logs into templates, which are then clustered based on events they correspond to, and the log clusters are processed to infer the sequencing relationships among events. 
The resulting graph of events is then enriched by weighting each arc with the average time passing between the logging of the source event and that of the target event.
The topology graph and the time-weighted control flow graphs constitute the baseline model against which to compare newly logged events. 
This is done by mapping the events logged by each service to their corresponding template and by checking whether its actual behaviour deviates from the expected one.
A service suffers from functional anomalies whenever it does not log an event that it should have been logging, or when unexpected logs occur.
A service instead suffers from performance anomalies when it logs an expected event significantly later than when expected. 
%If any anomaly is detected in a service, the service topology graph is used to look for anomalies also in the services depending on the anomalous one, where the anomaly may have propagated.

\subsection{Distributed Tracing-based Anomaly Detection Techniques}
\label{sec:detection:distributed-tracing}
% distributed tracing
Online anomaly detection is enacted also by instrumenting the target applications to feature additional functionalities, with the most common being distributed tracing.
Distributed tracing is indeed at the basis of various techniques for anomaly detection, combined with supervised or unsupervised machine learning,\footnote{Whilst unsupervised learning algorithms learn patterns directly from data, supervised learning algorithms rely on data to be labelled, \viz training examples associating given input data with desired output values \cite{Mitchell1997_MachineLearning}.} or with trace comparison techniques.

\subsubsection{Unsupervised Learning}
\label{sec:detection:distributed-tracing:unsupervised-learning}
TraceAnomaly \cite{16_Liu2020_TraceAnomaly} and Nedelkoski \etal \cite{35_Nedelkoski2019_AnomalyDetectionTracingData} collect traces in training runs of a multi-service application and use them to train unsupervised neural networks.
They then enact online anomaly detection by exploiting the trained neural networks to process the traces generated by the application while it is running in production.
% distributed tracing + neural network
In particular, TraceAnomaly \cite{16_Liu2020_TraceAnomaly} trains a deep Bayesian neural network with posterior flow \cite{Neal1996_BayesianNetwork,Rezende2015_NormalizingFlows}, which enables associating monitored traces with a likelihood to be normal, \viz not affected by performance anomalies.
It also stores all seen service call paths, \viz all sequences of service interactions observed in the available traces. 
TraceAnomaly \cite{16_Liu2020_TraceAnomaly} then enacts online anomaly detection in two steps.
It first checks whether a newly produced trace contains previously unseen call paths.
If this is the case, TraceAnomaly \cite{16_Liu2020_TraceAnomaly} decides whether to consider the unseen call paths as functional anomalies for the application, based on whether they appear in a whitelist manually filled by the application operator to exclude false alarms (\eg unseen call paths corresponding to new interactions occurring after service upgrades).
If there is no functional anomaly, the trace is passed to the deep Bayesian neural network, which associates the monitored trace with its likelihood to be normal.
If the likelihood is below a given threshold, the trace denotes a performance anomaly for the application.

Nedelkoski \etal \cite{35_Nedelkoski2019_AnomalyDetectionTracingData} train a multi-modal long-short term memory neural network \cite{Goodfellow2016_DeepLearningBook} modelling the normal behaviour of a multi-service application, based on the traces collected in failure-free runs of the application.
The long-short term memory neural network is multi-modal in the sense that it considers both the types of events logged in a trace and the response times of application services.
It is obtained by combining two single-modal networks, one trained to predict the probability distribution of the events following an event in a trace, and the other trained to predict the probability distribution of the response times for a service invocation.
The trained network is then used in online anomaly detection to predict the most probable events following an event and the most probable response times for a service invocation. 
If the event following another, or if the response time of a service are not amongst the predicted ones, the application is considered to suffer from a functional or performance anomaly, respectively.

% distributed tracing + machine learning
Jin \etal \cite{22_Jin2020_AnomalyDetectionMicroservicesRPCA} allows detecting performance anomalies in multi-service applications based on the offline analysis of the traces collected through distributed tracing, by also considering the performance metrics of application services that can be directly collected from their runtime environment (\eg CPU and memory consumption).
Jin \etal \cite{22_Jin2020_AnomalyDetectionMicroservicesRPCA} first enact principal component analysis on logged traces to determine the services that may be involved in anomalous interactions.
This is done by deriving a matrix-based representation of the logged traces, by reducing the obtained matrix to its principal components, and by applying a function linearly combining such principal components to elicit the services involved in anomalous traces. 
The performance metrics collected for such services are then processed with unsupervised learning algorithms for anomaly detection, \viz isolation forest, one-class support vector machine, local outlier factor, or 3$\sigma$ \cite{Chandola2009_SurveyAnomalyDetection}. 
Such algorithms are used to determine anomalous values for such metrics, which are used to assign an anomaly score to the services involved in anomalous traces.
Jin \etal \cite{22_Jin2020_AnomalyDetectionMicroservicesRPCA} then return the list of services affected by performance anomalies, \viz the services whose anomaly score is higher than a given threshold, ranked by anomaly score.

\subsubsection{Supervised Learning}
Seer \cite{17_Gan2019_Seer}, Nedelkoski \etal \cite{34_Nedelkoski2019_AnomalyDetectionDistributedTracing}, and Bogatinovski \etal \cite{36_Bogatinovski2020_SelfSupervisedAnomalyDetection} also enact a combination of distributed tracing and deep learning, however following a supervised learning approach.
They indeed collect traces in training runs of a multi-service application, labelled to distinguish normal runs from those where anomalies are known to have occurred on specific services.
The collected traces are used to train neural networks, which are then used to enact online anomaly detection by processing the traces newly generated by the application.
% distributed tracing + neural network
In particular, Seer \cite{17_Gan2019_Seer} trains a deep neural network consisting of convolutional layers followed by long-short term memory layers \cite{Goodfellow2016_DeepLearningBook}. 
Input and output neurons of the network corresponds to application services.
Input neurons are used to pass the KPIs of each service, \viz latency and outstanding requests, logged by the service itself in distributed traces, and resource consumption, obtained by interacting with the nodes where the services are running. 
Output neurons instead indicate which services are affected by performance anomalies.
At runtime, Seer \cite{17_Gan2019_Seer} continuously feeds the trained deep neural network with the stream of monitored traces, and the network consumes such traces to determine whether some application service is affected by a performance anomaly. 
If this is the case, Seer \cite{17_Gan2019_Seer} interacts with the runtime of the node hosting the anomalous service to determine which computing resources are being saturated.
Seer \cite{17_Gan2019_Seer} then notifies the system manager about the resource saturation happening on the node, to allow the manager to mitigate the effects of the detected performance degradation, \eg by providing the node with additional computing resources.

Nedelkoski \etal \cite{34_Nedelkoski2019_AnomalyDetectionDistributedTracing} train two neural networks, \viz a variational autoencoder \cite{Kingma2015_VariationalAutoencoder} modelling the normal behaviour of the application and a convolutional neural network \cite{Goodfellow2016_DeepLearningBook} recognising the type of failure affecting a service when the latter is affected by a performance anomaly. 
The variational autoencoder is learned from the distributed traces obtained in normal runs of the application, so that the autoencoder learns to encode non-anomalous traces and to suitably reconstruct them from their encoding.
If applied on an anomalous trace, instead, the autoencoder reconstruct them with significant errors, hence enabling to detect anomalies.
The convolutional neural network is instead trained on the traces collected in runs of the application where specific failures were injected in specific services, so that the network can recognise which failures caused a performance anomaly in a service.
In online anomaly detection, the trained networks are used in a pipeline.
The outputs of the autoencoder are passed to a post-processor to exclude false positives, \viz performance anomalies affecting a service due to temporary reasons (\eg overload of a service). 
%The post-processor works with a tolerance threshold, which can be customised to tailor it to the multi-service application under consideration and its runtime environment. 
If a service is considered anomalous by the post-processor as well, the anomalous trace is passed to the convolutional network to detect the type of anomaly affecting the service. 

Bogatinovski \etal \cite{36_Bogatinovski2020_SelfSupervisedAnomalyDetection} work under a different assumption if compared with the other anomaly detection techniques discussed above, \viz it assumes that the logging of an event in a particular position in a trace is conditioned both by those logged earlier and by those appearing afterwards.
Bogatinovski \etal \cite{36_Bogatinovski2020_SelfSupervisedAnomalyDetection} indeed train a self-supervised encoder-decoder neural network \cite{Goodfellow2016_DeepLearningBook} capable of predicting the logging of an event in a given \enquote{masked} position of a trace based on the context given by its neighbour events. 
In particular, the trained neural network provides a probability distribution of the possible events appearing in any masked position of an input trace, given the context of the events appearing before and after the masked event to be predicted. 
The trained neural network is then used to enact anomaly detection over novel traces, \viz to obtain a sorted list of events that should have been logged in each position of the trace. 
The lists generated by the network are then analysed by a post-processor, which considers a truly logged event as anomalous if it is not amongst the events predicted to appear in the corresponding position. 
The post-processor then computes an anomaly score for the trace (\viz number of anomalous events divided by trace length).
If the anomaly score is beyond an user-specified threshold, the trace is considered to witness a functional anomaly affecting the application.

% distributed tracing + machine learning
Finally, MEPFL \cite{02_Zhou2019_LatentErrorPrediction} provides another technique applying supervised machine learning to perform online anomaly detection on application instrumented to feature distributed tracing.
MEPFL \cite{02_Zhou2019_LatentErrorPrediction} can recognize functional anomalies by exploiting multiple supervised learning algorithms, \viz k-nearest neighbors \cite{Altman1992_NearestNeighbor}, random forests \cite{Breiman2001_RandomForests}, and multi-layer perceptron \cite{Goodfellow2016_DeepLearningBook}.
MEPFL \cite{02_Zhou2019_LatentErrorPrediction} takes as input a multi-service application and a set of automated test cases simulating user requests to load the application and to produce traces.
The traces must include information on the configuration of a service (\viz memory and CPU limits, and volume availability), as well as on the status, resource consumption, and service interactions of its instances.
This information is fed to one of the supported supervised learning algorithms, which trains classifiers determining whether a trace includes anomalies, which services are affected by such anomalies, and which types of failures caused such services to experience anomalies. 
The training of the classifiers is based on traces collected in training runs of an application, where the application is loaded with the input test cases, and where the application is run both in normal conditions and by injecting failures in its services to observe how the traces change when such type of failures affect such services.

\subsubsection{Trace Comparison}
Trace comparison is another possible technique to enact online anomaly detection on multi-service applications instrumented to feature distributed tracing.
This is the technique followed by Meng \etal \cite{23_Meng2021_DetectingAnomaliesMicroservices}, Wang \etal \cite{26_Wang2020_WorkflowAwareFaultDiagnosis}, and Chen \etal \cite{39_Chen2020_MatrixSketchBasedAnomalyDetection}, which all rely on collecting the traces that can possibly occur in a multi-service application and by then checking whether newly collected traces are similar to the collected ones.

% distributed tracing + pca
Meng \etal \cite{23_Meng2021_DetectingAnomaliesMicroservices} and Wang \etal \cite{26_Wang2020_WorkflowAwareFaultDiagnosis} enable detecting functional anomalies by collecting traces while testing an application in a pre-production environment and by building a set of representative call trees from collected traces.
Such trees model the service invocation chains that can possibly appear in normal conditions, as they have been observed in the collected traces. %, by including each chain only once, even if the corresponding sequence of service calls was repeated multiple times (\eg because of loops or recursion in the callee). 
Newly monitored traces are then reduced to their call trees, whose tree-edit distance from each representative call tree is computed with the RTDM algorithm \cite{Reis2004_RTDM}.
If the minimum among all the computed distances is higher than a given threshold, the trace is considered as anomalous, and the service from which tree edits start is identified as the service suffering from a functional anomaly.
Meng \etal \cite{23_Meng2021_DetectingAnomaliesMicroservices} and Wang \etal \cite{26_Wang2020_WorkflowAwareFaultDiagnosis} also enable detecting performance anomalies based on the services' response time.
The response times in the traces collected in the testing phase are modelled as a matrix, whose elements represent the response time for a service in each of the collected traces.
Principal component analysis \cite{Jolliffe2011_PCA} is then enacted to reduce the matrix to its principal components, which are then used to define a function determining whether the response time of a service in newly monitored traces is anomalous.
At the same time, as both Meng \etal \cite{23_Meng2021_DetectingAnomaliesMicroservices} and Wang \etal \cite{20_Wang2020_RootCauseMetricLocation} explicitly notice, the computational complexity of the enacted trace comparisons makes them better suited to enact offline anomaly detection, as they would be too time consuming to perform online anomaly detection on medium-/large-scale multi-service applications. 

% distributed tracing + matrix sketch (dimensionality reduction as in pca, but different technique)
Chen \etal \cite{39_Chen2020_MatrixSketchBasedAnomalyDetection} instead uses matrix sketching to compare traces and detect anomalies.
More precisely, Chen \etal \cite{39_Chen2020_MatrixSketchBasedAnomalyDetection} adapt an existing matrix sketching algorithm \cite{Huang2015_MatrixSketching} to detect response time anomalies in services in monitored traces. 
This is done by maintaining up-to-date a \enquote{sketch} of the high-dimensional data corresponding to the historically monitored response times of the services forming an application.
The sketch is a limited set of orthogonal basis vectors that can linearly reconstruct the space containing the non-anomalous response times for each service in the formerly monitored traces. 
The response time for a service in a newly monitored trace is non-anomalous if it lies within or close enough to the space defined by the sketch vector, given a tolerance threshold.
If this is the case, the sketch vectors are updated.
Otherwise, the corresponding service is considered to experience a performance anomaly.

\subsection{Monitoring-based Anomaly Detection Techniques}
\label{sec:detection:monitoring}
Online anomaly detection in multi-service applications can also be enacted by installing agents to monitor KPIs on their services, and by then processing such KPIs to detect anomalies.
In this perspective, a basic solution is comparing the KPIs monitored on the service acting as the application frontend against the application's SLOs: in case of SLO violations, the whole application is considered to suffer from a performance anomaly.
For detecting anomalies at a finer granularity, \viz going from application-level anomalies to service-level anomalies, the most common technique is to exploit unsupervised or supervised machine learning algorithms to process monitored KPIs and detect whether any anomaly is occurring in any application service.
Another possibility is to exploit self-adaptive heartbeat protocols to detect anomalous services, rather than processing their KPIs. 
We hereafter discuss all such techniques by starting from those based on unsupervised/supervised learning, which are then followed by those based on SLO checks and heartbeating.

\subsubsection{Unsupervised Learning}
\label{sec:detection:monitoring:unsupervised-learning}
Online anomaly detection can be enacted by exploiting unsupervised learning algorithms to process the KPIs monitored on application services in failure-free runs of the application and learn a baseline modelling of their behaviour.
The baseline model is then used online to detect whether the newly monitored KPIs on a service diverge from the baseline model, hence indicating that the service suffers from some anomaly.
This technique is adopted by Gulenko \etal \cite{37_Gulenko2018_DetectingAnomalousBehaviourBlackBoxServices}, MicroRCA \cite{04_Wu2020_MicroRCA}, Wu \etal \cite{09_Wu2020_MicroserviceDiagnosisDeepLearning}, LOUD \cite{13_Mariani2018_LocalizingFaultsCloudSystems}, and DLA \cite{14_Samir2019_DLA}, which train the baseline model by considering the KPIs monitored in normal runs of the application, \viz assuming that no anomaly occurred in such runs.

Gulenko \etal \cite{37_Gulenko2018_DetectingAnomalousBehaviourBlackBoxServices} and LOUD \cite{13_Mariani2018_LocalizingFaultsCloudSystems} are \enquote{KPI-agnostic}, meaning that they can work on any set of KPIs for the services forming an application, monitored with probes installed in the hosts where the application services are run.  
% BIRCH clustering for baseline clusters 
In particular, Gulenko \etal \cite{37_Gulenko2018_DetectingAnomalousBehaviourBlackBoxServices} enable detecting performance anomalies in multi-service applications by building different baseline models for different services.
The KPIs monitored on a service at a given time are modelled as a vector, and the BIRCH \cite{Zhang1996_BIRCH} online clustering algorithm is used in an initial training phase to cluster the vectors corresponding to the monitored KPIs.
At the end of the training phase, the centroids and centroid radii of the obtained clusters determine the baseline subspaces where the vectors of monitored KPIs should pertain to be classified as \enquote{normal}. 
This enables the online processing of newly monitored KPIs for each service, by simply checking whether the corresponding vector falls within any of the baseline subspaces, \viz whether there exists a cluster whose centroid is less distant from the vector of monitored KPIs than the corresponding centroid radius. 
If this is not the case, the corresponding service is considered as affected by a performance anomaly.

LOUD \cite{13_Mariani2018_LocalizingFaultsCloudSystems} instead trains a baseline model for the whole application by pairing each monitored KPI with the corresponding service.
The KPIs monitored in the offline runs of the application are passed to the IBM ITOA-PI \cite{IBM_ITOA_PI} to build a KPI baseline model and a causality graph.
The baseline model provides information about the average values and the acceptable variations over time for a KPI.
Each node in the causality graph correspond to a service's KPI, and the edges indicate the causal relationships among KPIs, with weights indicating the probability to which changes in the source KPI can cause changes in the target KPI (computed with the Granger causality test \cite{Arnold2007_GrangerMethodsTemporalCausalModelling}).
During the online phase, LOUD \cite{13_Mariani2018_LocalizingFaultsCloudSystems} again exploits the IBM ITOA-PI to detect performance anomalies: a KPI and the corresponding service are reported as anomalous if the monitored value for such KPI is outside of the acceptable variations coded in the baseline model, or if the causal relationships represented in the causality graph are not respected (\eg when a KPI significantly changed its value, but those that should have changed as well remained unchanged).

Differently from the above discussed, \enquote{KPI-agnostic} techniques \cite{37_Gulenko2018_DetectingAnomalousBehaviourBlackBoxServices,13_Mariani2018_LocalizingFaultsCloudSystems}, MicroRCA \cite{04_Wu2020_MicroRCA}, Wu \etal \cite{09_Wu2020_MicroserviceDiagnosisDeepLearning}, and DLA \cite{14_Samir2019_DLA} focus on given KPIs.
MicroRCA \cite{04_Wu2020_MicroRCA} and Wu \etal \cite{09_Wu2020_MicroserviceDiagnosisDeepLearning} consider the response time and resource consumption of application services, which they monitor by requiring the Kubernetes (k8s) deployment of multi-service applications to be instrumented as service meshes featuring Istio \cite{Istio} and Prometheus \cite{Prometheus}.
The latter are used to collect KPIs from the application services. %, their containers, and the nodes used to host them. 
MicroRCA \cite{04_Wu2020_MicroRCA} and Wu \etal \cite{09_Wu2020_MicroserviceDiagnosisDeepLearning} then follow an unsupervised learning approach similar to that proposed by Gulenko \etal \cite{37_Gulenko2018_DetectingAnomalousBehaviourBlackBoxServices}.
The KPIs monitored on each service %and on its container and host 
are modelled as a vector, and the BIRCH algorithm \cite{Zhang1996_BIRCH} %online clustering algorithm 
is used to cluster the vectors corresponding to KPIs monitored in normal runs of the application, \viz assuming that no anomaly occurred in such runs.
The obtained clusters determine the baseline subspaces where the vectors of monitored KPIs should pertain to be classified as \enquote{normal}. 
This enables the online processing of newly monitored KPIs for each service, by simply checking whether the corresponding vector falls within any of the baseline subspaces. 
If this is not the case, the service is considered as affected by a performance anomaly.

DLA \cite{14_Samir2019_DLA} instead focuses on the response time of each application service, whilst also considering its possible fluctuations due to increase/decrease of end users' transactions.
It also focuses on containerised multi-service applications, whose deployment in k8s is expected to be provided by the application operator. 
DLA \cite{14_Samir2019_DLA} installs a monitoring agent in the VMs used to deploy the application.
The agents collect the response time of the containerised services running in each VM, by associating each monitored response time with a timestamp and with the corresponding user transaction. 
This information is used to learn whether/how the response time varies based on the number of user transactions occurring simultaneously, based on the Spearman’s rank correlation coefficient \cite{Sheskin2011_ParametricNonparametricStatisticalProcedures}.
The obtained correlation is then used to check whether a variation in a newly monitored response time corresponds to an expected fluctuation due to an increase/decrease of users’ transactions, or whether it actually denotes a performance anomaly affecting a service.

The above discussed techniques train a baseline model of the application behaviour in an offline learning step, and then employ the trained model to detect anomalies while the application is running. 
This technique works under the assumption that the conditions under which an application is run do not change over time (\eg no new application is deployed on the same VMs) or that what monitored during the training phase is enough to model all possible situations for the application \cite{Xu2017_AdaptiveServiceAPIMonitoring}. 
%DLA \cite{14_Samir2019_DLA} actually recognises this possible issue, also showing how the offline training can be periodically repeated to keep the baseline models updated to the new conditions under which an application is running.
CloudRanger \cite{29_Wang2018_CloudRanger} tackles the problem from a different perspective, with the ultimate goal of continuously adapting the anomaly detection system to follow the evolution of the application and of the context where it runs. 
CloudRanger \cite{29_Wang2018_CloudRanger} indeed enacts continuous learning on the frontend service of an application, to learn which is its \enquote{expected} response times in the current conditions, \viz based on the monitored response time, load, and throughput.
This is done with the initial, offline training of a baseline model on historically monitored data for the application frontend. %, \viz estimating the expected response time for given load and throughput. 
When the application is running, CloudRanger \cite{29_Wang2018_CloudRanger} applies polynomial regression to monitored KPIs to update the baseline model to reflect the application behaviour in the dynamically changing runtime conditions.
Online anomaly detection is then enacted by checking whether the response time of the application frontend is significantly different from the expected one, \viz if the difference between the expected and monitored response times is higher than a given threshold.

% alternative technique: anomaly prediction (value + time)
Finally, Hora \cite{27_Pitakrat2018_HoraExtended,28_Pitakrat2016_Hora} tackles the problem of online anomaly detection from yet another perspective, different from all those discussed above.
It indeed combines architectural models with statistical analysis techniques to preemptively determine the occurrence of performance anomalies in a multi-service application. 
Hora exploits SLAstic \cite{vanHoorn2014_SLAstic} to automatically extract a representation of the architectural entities in a multi-service application (\viz application services and nodes used to host them) and the degree to which a component depends on another.
This information is used to create Bayesian networks \cite{Neal1996_BayesianNetwork} modelling the possible propagation of performance anomalies among the services forming an application.
The online anomaly detection is then enacted by associating each application service with a detector, which monitors given KPIs for such service, and which analyses such monitored metrics to detect whether the currently monitored KPIs denote a performance anomaly. 
Since monitored metrics are time series data, detectors exploit autoregressive integrated moving average \cite{Shumway2017_ARIMA} to determine whether performance anomalies affect monitored services. 
Whenever a performance anomaly is detected on a service, Hora \cite{27_Pitakrat2018_HoraExtended,28_Pitakrat2016_Hora} enacts Bayesian inference \cite{Neal1996_BayesianNetwork} to determine whether the anomaly propagates to other services. 
It then returns the set of application services most probably being affected by a performance anomaly.

\subsubsection{Supervised Learning}
ADS \cite{15_Du2018_AnomalyDetectionContainerMicroservices} and PreMiSE \cite{33_Mariani2020_PredictingFailures} also enact online anomaly detection by training a baseline model on monitored KPIs.
They however enact supervised learning (rather than unsupervised learning), by also injecting specific failures in specific services during the training phase.
They indeed label monitored KPIs to denote whether they correspond to normal runs or to runs where specific failures were injected on specific services.
The monitoring is enacted by installing monitoring agents in the VMs used to deploy an application, even if the two techniques consider different types of application deployments and KPIs.
As a result, ADS \cite{15_Du2018_AnomalyDetectionContainerMicroservices} and PreMiSE \cite{33_Mariani2020_PredictingFailures} can model both the normal behaviour of the application services and their anomalous behaviour when failures similar to those injected occur.


% baseline model of microservices, including possible failures
ADS \cite{15_Du2018_AnomalyDetectionContainerMicroservices} enables detecting performance anomalies in containerised multi-service applications, given their k8s deployment and modules for injecting failures in their services.
ADS \cite{15_Du2018_AnomalyDetectionContainerMicroservices} considers a predefined set of KPIs for each containerised application service, \viz CPU, memory, and network consumption.
In the training phase, the application is loaded with a workload generator, which is provided by the application operator to simulate end-user requests. 
Multiple runs of the application are simulated, both in failure-free conditions and by exploiting the fault injection modules to inject failures in the application services. 
The monitored KPIs are labelled as pertaining to \enquote{normal} or \enquote{anomalous} runs of the application, in such a way that it can be processed to train a classifier with a supervised machine learning algorithm (\viz nearest neighbour \cite{Altman1992_NearestNeighbor}, random forest \cite{Breiman2001_RandomForests}, na\"ive Bayes \cite{John1995_NaiveBayes}, or support vector machines \cite{Scholkopf1999_SupportVectorLearning}), assuming that no more than one anomaly occurs at the same time.
The trained classifier provides the baseline model for normal/anomalous behaviour of the application, and it is used for classifying application services as experiencing some known performance anomaly based on their newly monitored KPIs.

PreMiSE \cite{33_Mariani2020_PredictingFailures} trains a baseline model for detecting anomalies in multi-service applications, assuming that each service runs in a different VM. 
Monitored KPIs are treated as time series: in an intial, offline learning phase, PreMiSE \cite{33_Mariani2020_PredictingFailures} trains a baseline modelling of the normal, failure-free execution of the application services. 
The baseline model captures temporal relationships in the time series data corresponding to monitored KPIs to model trends and seasonalities, and it applies Granger causality tests \cite{Arnold2007_GrangerMethodsTemporalCausalModelling} to determine whether the correlation among two KPIs is such that one can predict the evolution of the other.
PreMiSE \cite{33_Mariani2020_PredictingFailures} also trains a signature model representing the anomalous behaviour of the application in presence of failures.
This is done by injecting predefined failures (\viz packet loss/corruption, increased network latency, memory leak, and CPU hog) in the VMs running the application services and by monitoring the corresponding changes in KPIs.
The monitored KPIs are then used to train a classifier for detecting the occurrence of the anomalies corresponding to the injected failures, with one out of six supported supervised machine learning algorithms. 
The trained baseline model is then compared against newly monitored KPIs to enact online anomaly detection. 
Once an anomaly is detected, the signature model is used to classify the application service suffering from the detected anomaly and the failure it is suffering from. 

% basic: SLO monitoring of the fronted
\subsubsection{SLO Check}
\label{sec:detection:monitoring:slo-check}
CauseInfer \cite{40_Chen2020_CauseInfer,48_Chen2014_CauseInfer} and Microscope \cite{06_Lin2018_Microscope,19_Guan2019_MicroscopeDemo} enact online anomaly detection in multi-service applications by monitoring KPIs of the application frontend and comparing such KPIs with the SLOs of the application.
In particular, CauseInfer \cite{40_Chen2020_CauseInfer,48_Chen2014_CauseInfer} exploits a monitoring agent to continuously monitor the response time of the service acting as frontend for an application.
Microscope \cite{06_Lin2018_Microscope,19_Guan2019_MicroscopeDemo} instead relies on the k8s deployment of an application to include monitoring agents (like Prometheus \cite{Prometheus}), throughout which application services expose their response time.
CauseInfer \cite{40_Chen2020_CauseInfer,48_Chen2014_CauseInfer} and Microscope \cite{06_Lin2018_Microscope} and then compare the KPIs monitored on the frontend of an application with that declared in the SLOs of the application. 
Whenever the performance of the application frontend violates the application SLOs, CauseInfer \cite{40_Chen2020_CauseInfer,48_Chen2014_CauseInfer} and Microscope \cite{06_Lin2018_Microscope,19_Guan2019_MicroscopeDemo} consider the whole application as affected by a performance anomaly.

% CauseInfer \cite{40_Chen2020_CauseInfer,48_Chen2014_CauseInfer} enacts online anomaly detection in multi-service applications by monitoring KPIs of the application frontend and comparing such KPIs with the SLOs of the application.
% In particular, CauseInfer \cite{40_Chen2020_CauseInfer,48_Chen2014_CauseInfer} exploits a monitoring agent to continuously monitor the response time of the service acting as frontend for an application, and it compares such response time with that declared in the SLOs of the application.
% Whenever the performance of the application frontend violates the application SLOs, CauseInfer \cite{40_Chen2020_CauseInfer,48_Chen2014_CauseInfer} considers the whole application as affected by a performance anomaly.

A similar technique is adopted by $\epsilon$-diagnosis \cite{49_Shan2019_epsilonDiagnosis}, which also monitors SLOs on the frontend service of an application to enact online performance anomaly detection.
$\epsilon$-diagnosis \cite{49_Shan2019_epsilonDiagnosis} however focuses on the tail latency of the application frontend, determined in small time windows (\eg one minute or one second), rather than on bigger ones as typically done. 
Given the k8s deployment of the target application, the $\epsilon$-diagnosis system \cite{49_Shan2019_epsilonDiagnosis} enacts online anomaly detection as follows: it partitions the stream of data into small-sized windows, it computes the tail latency for each window, and it then compares the tail latency with the threshold declared in the application SLOs.
If the tail latency for a time window is higher than the given threshold, the whole application is considered to suffer from a performance anomaly in such a time window.

\subsubsection{Heartbeating}
% heartbeat 
A different, monitoring-based technique to detect anomalous services in multi-service applications is heartbeating.
The latter requires installing a monitorer sending heartbeat messages to the application services, which must reply to such messages within a given timeframe to be considered as fully working.
M-MFSA-HDA \cite{24_Zang2018_FaultDiagnosisMicroservices} adopts this technique, by periodically sending heartbeat messages to the services forming an application.
If any of such services does not reply within a given timeframe, it is considered to suffer from a functional anomaly. %, most probably due to the fact that it failed.
In addition, to keep a suitable trade-off between the impact of the heartbeat protocol on application performances and the failure detection performances of the heartbeat detection system itself, M-MFSA-HDA \cite{24_Zang2018_FaultDiagnosisMicroservices} self-adaptively changes the heartbeat rate by also monitoring the CPU consumption of the nodes hosting application services, the network load, and the application workload. 

\subsection{Discussion}
\label{sec:detection:discussion}

%\begin{table}[tp]
\caption{Classification of anomaly detection techniques, based on their \textit{class} (\viz \textsf{L}, \textsf{DT}, and \textsf{M} for log-based, distributed tracing-based, and monitoring-based techniques, respectively), the applied \textit{method}, the \textit{type} (\viz \textsf{F} and \textsf{P} for functional and performance anomalies, respectively) and \textit{granularity} of detected anomalies (\viz \textsf{A} and \textsf{S} for application-level and service-level anomalies, respectively), and the \textit{artifacts} they need to run.}
\label{tab:detection}
\begin{footnotesize}
\begin{tabular}{%
    >{\centering\arraybackslash}m{.2\textwidth}%
    @{\ }
    >{\centering\arraybackslash}m{.06\textwidth}%
    @{\ }
    >{\centering\arraybackslash}m{.21\textwidth}%
    @{\ }
    >{\centering\arraybackslash}m{.05\textwidth}%
    @{\,}
    >{\centering\arraybackslash}m{.11\textwidth}%
    @{\ }
    >{\centering\arraybackslash}m{.32\textwidth}%
}
    \thickline
    %study
        & 
        &
        & \multicolumn{2}{c}{\textbf{Anomaly}}
        & 
        \\ 
    %study
        & \textbf{Class} 
        & \textbf{Detection Method}
        & \textbf{Type}
        & \textbf{Granularity} 
        & \textbf{Needed Artifacts}
        \\ 
    \thickline
    OASIS \cite{38_Nandi2016_AnomalyDetectionControlFlowGraphMining}
        & \textsf{L} 
        & \cellcolor[HTML]{FDF3ED} unsupervised learning
        %& online
        & \textsf{F}
        & \textsf{A}
        & training logs, runtime logs
        \\
    Jia \etal \cite{31_Jia2017_AnomalyDiagnosisHybridGraphModel}
        & \textsf{L} 
        & \cellcolor[HTML]{FDF3ED} unsupervised learning
        %& online
        & \textsf{F}, \textsf{P}
        & \textsf{S}
        & training logs, runtime logs
        \\
    LogSed \cite{32_Jia2017_LogSed}
        & \textsf{L} 
        & \cellcolor[HTML]{FDF3ED} unsupervised learning
        %& online
        & \textsf{F}, \textsf{P}
        & \textsf{S}
        & training logs, runtime logs
        \\
        \thinline
    TraceAnomaly \cite{16_Liu2020_TraceAnomaly}
        & \textsf{DT} 
        & \cellcolor[HTML]{FDF3ED} unsupervised learning
        %& online
        & \textsf{F}, \textsf{P}
        & \textsf{A}
        & training traces, runtime traces
        \\
    Nedelkoski \etal \cite{35_Nedelkoski2019_AnomalyDetectionTracingData}
        & \textsf{DT}  %(with Zipkin
        & \cellcolor[HTML]{FDF3ED} unsupervised learning 
        %& online
        & \textsf{F}, \textsf{P}
        & \textsf{A}
        & training traces, runtime traces
        \\
    Jin \etal \cite{22_Jin2020_AnomalyDetectionMicroservicesRPCA}
        & \textsf{DT} 
        & \cellcolor[HTML]{FDF3ED} unsupervised learning 
        %& offline % the authors explicitly state that their approach works well only offline
        & \textsf{P}
        & \textsf{S}
        & reference traces, runtime traces
        \\
    
    Seer \cite{17_Gan2019_Seer}
        & \textsf{DT} 
        & \cellcolor[HTML]{D9FDE5} supervised learning
        %& online
        & \textsf{P}
        & \textsf{S}
        & training traces, runtime traces
        \\
    
    Nedelkoski \etal \cite{34_Nedelkoski2019_AnomalyDetectionDistributedTracing}
        & \textsf{DT}  %(with Zipkin
        & \cellcolor[HTML]{D9FDE5} supervised learning
        %& online
        & \textsf{P}
        & \textsf{S}
        & training traces, runtime traces
        \\
    
    Bogatinovski \etal \cite{36_Bogatinovski2020_SelfSupervisedAnomalyDetection}
        & \textsf{DT} 
        & \cellcolor[HTML]{D9FDE5} supervised learning
        & \textsf{P}
        & \textsf{A}
        & training traces, runtime traces
        \\
    
    MEPFL \cite{02_Zhou2019_LatentErrorPrediction}
        & \textsf{DT} 
        & \cellcolor[HTML]{D9FDE5} supervised learning 
        %& online               
        & \textsf{F}
        & \textsf{S}
        & target application deployment, automated test cases
        \\
    Meng \etal \cite{23_Meng2021_DetectingAnomaliesMicroservices}
        & \textsf{DT}  %(with Jaeger, Zipkin
        & \cellcolor[HTML]{dce3ac} trace comparison
        %& offline % in TTV, authors state that the proposed approach is too costly to enact online detection on medium/large-scale multi-service apps
        & \textsf{F}, \textsf{P}
        & \textsf{S}
        & reference traces, runtime traces
        \\
    Wang \etal \cite{26_Wang2020_WorkflowAwareFaultDiagnosis}
        & \textsf{DT}  %(with OpenTracing
        & \cellcolor[HTML]{dce3ac} trace comparison
        %& offline % in TTV, authors state that the proposed approach is too costly to enact online detection on medium/large-scale multi-service apps
        & \textsf{F}, \textsf{P}
        & \textsf{S}
        & reference traces, runtime traces
        \\
    Chen \etal \cite{39_Chen2020_MatrixSketchBasedAnomalyDetection}
        & \textsf{DT} 
        & \cellcolor[HTML]{dce3ac} trace comparison
        %& online
        & \textsf{P}
        & \textsf{S}
        & reference traces, runtime traces
        \\
        \thinline
    Gulenko \etal \cite{37_Gulenko2018_DetectingAnomalousBehaviourBlackBoxServices}
        & \textsf{M} 
        & \cellcolor[HTML]{FDF3ED} unsupervised learning
        %& online
        & \textsf{P}
        & \textsf{S}
        & target application deployment, workload generator
        \\
    LOUD \cite{13_Mariani2018_LocalizingFaultsCloudSystems}
        & \textsf{M}  
        & \cellcolor[HTML]{FDF3ED} unsupervised learning
        %& online
        & \textsf{P}
        & \textsf{S}
        & target application deployment, workload generator
        \\
    MicroRCA \cite{04_Wu2020_MicroRCA} 
        & \textsf{M} 
        & \cellcolor[HTML]{FDF3ED} unsupervised learning
        %& online
        & \textsf{P}
        & \textsf{S}
        & k8s deployment, workload generator
        \\ 
    
    Wu \etal \cite{09_Wu2020_MicroserviceDiagnosisDeepLearning}
        & \textsf{M} 
        & \cellcolor[HTML]{FDF3ED} unsupervised learning
        %& online
        & \textsf{P}
        & \textsf{S}
        & k8s deployment, workload generator
        \\
    
    DLA \cite{14_Samir2019_DLA}
        & \textsf{M} 
        & \cellcolor[HTML]{FDF3ED} unsupervised learning
        %& online
        & \textsf{P}
        & \textsf{S}
        & k8s deployment % HHMM model - not for anomaly detection!
        \\
    
    CloudRanger \cite{29_Wang2018_CloudRanger}
        & \textsf{M} 
        & \cellcolor[HTML]{FDF3ED} unsupervised learning
        %& online
        & \textsf{P}
        & \textsf{A}
        & target application deployment, historically monitored KPIs
        \\
    
    Hora \cite{27_Pitakrat2018_HoraExtended,28_Pitakrat2016_Hora}
        & \textsf{M} 
        & \cellcolor[HTML]{FDF3ED} unsupervised learning
        %& online
        & \textsf{P}
        & \textsf{S}
        & target application deployment
        \\
    ADS \cite{15_Du2018_AnomalyDetectionContainerMicroservices}
        & \textsf{M} 
        & \cellcolor[HTML]{D9FDE5} supervised learning
        & \textsf{P}
        & \textsf{S}
        & k8s deployment, failure injectors, workload generator
        \\
    PreMiSE \cite{33_Mariani2020_PredictingFailures}
        & \textsf{M} 
        & \cellcolor[HTML]{D9FDE5} supervised learning
        %& online
        & \textsf{F}, \textsf{P}
        & \textsf{S}
        & target application deployment, workload generator % with one service per VM
        \\
    
    CauseInfer \cite{40_Chen2020_CauseInfer,48_Chen2014_CauseInfer}
        & \textsf{M} 
        & \cellcolor[HTML]{d1bfb3} SLO check
        %& online
        & \textsf{P}
        & \textsf{A}
        & target application deployment, SLOs
        \\
    
    MicroScope \cite{06_Lin2018_Microscope,19_Guan2019_MicroscopeDemo}
        & \textsf{M} %\textsf{M} 
        & \cellcolor[HTML]{d1bfb3} SLO check
        %& online
        & \textsf{P}
        & \textsf{A}
        & k8s deployment, SLOs
        \\
    
    $\epsilon$-diagnosis \cite{49_Shan2019_epsilonDiagnosis}
        & \textsf{M} 
        & \cellcolor[HTML]{d1bfb3} SLO check
        %& online
        & \textsf{P}
        & \textsf{A}
        & k8s deployment, SLOs
        \\
    M-MFSA- HDA \cite{24_Zang2018_FaultDiagnosisMicroservices}
        & \textsf{M} 
        & \cellcolor[HTML]{7BE0C5} heartbeating
        %& online
        & \textsf{F}
        & \textsf{S}
        & target application deployment
        \\
    \thickline
\end{tabular}
\end{footnotesize}
\end{table}
\begin{table}[tp]
\caption{Classification of anomaly detection techniques, based on their \textit{class} (\viz \textsf{L} for log-based techniques, \textsf{DT} for distributed tracing-based techniques, \textsf{M} for monitoring-based techniques), the applied \textit{method}, the \textit{type} (\viz \textsf{F} for functional anomalies, \textsf{P} for performance anomalies) and \textit{granularity} of detected anomalies (\viz \textsf{A} for application-level anomalies, \textsf{S} service-level anomalies), and the \textit{input} they need to run.}
\label{tab:detection}
\begin{tiny}
\begin{tabular}{%
    >{\centering\arraybackslash}m{.15\textwidth}%
    @{\ }
    >{\centering\arraybackslash}m{.06\textwidth}%
    @{\ }
    >{\centering\arraybackslash}m{.18\textwidth}%
    @{\ }
    >{\centering\arraybackslash}m{.05\textwidth}%
    @{\,}
    >{\centering\arraybackslash}m{.08\textwidth}%
    @{\ }
    >{\centering\arraybackslash}m{.4\textwidth}%
}
    \thickline
    %study
        & 
        &
        & \multicolumn{2}{c}{\textbf{Anomaly \ }}
        & 
        \\ 
    %study
    \textbf{Reference} 
        & \textbf{Class} 
        & \textbf{Method}
        & \textbf{Type}
        & \textbf{Gran.} 
        & \textbf{Needed Input}
        \\ 
    \thickline
    OASIS \cite{38_Nandi2016_AnomalyDetectionControlFlowGraphMining}
        & \textsf{L} 
        & \cellcolor[HTML]{E7F2F8} unsupervised learning
        %& online
        & \textsf{F}
        & \textsf{A}
        & previous and runtime logs
        \\
    Jia \etal \cite{31_Jia2017_AnomalyDiagnosisHybridGraphModel}
        & \textsf{L} 
        & \cellcolor[HTML]{E7F2F8} unsupervised learning
        %& online
        & \textsf{F}, \textsf{P}
        & \textsf{S}
        & previous and runtime logs
        \\
    LogSed \cite{32_Jia2017_LogSed}
        & \textsf{L} 
        & \cellcolor[HTML]{E7F2F8} unsupervised learning
        %& online
        & \textsf{F}, \textsf{P}
        & \textsf{S}
        & previous and runtime logs
        \\
       \thinline
    TraceAnomaly \cite{16_Liu2020_TraceAnomaly}
        & \textsf{DT} 
        & \cellcolor[HTML]{E7F2F8} unsupervised learning
        %& online
        & \textsf{F}, \textsf{P}
        & \textsf{A}
        & previous and runtime traces
        \\
    Nedelkoski \etal \cite{35_Nedelkoski2019_AnomalyDetectionTracingData}
        & \textsf{DT}  %(with Zipkin
        & \cellcolor[HTML]{E7F2F8} unsupervised learning 
        %& online
        & \textsf{F}, \textsf{P}
        & \textsf{A}
        & previous and runtime traces
        \\
    Jin \etal \cite{22_Jin2020_AnomalyDetectionMicroservicesRPCA}
        & \textsf{DT} 
        & \cellcolor[HTML]{E7F2F8} unsupervised learning 
        %& offline % the authors explicitly state that their approach works well only offline
        & \textsf{P}
        & \textsf{S}
        & previous and runtime traces
        \\
    
    Seer \cite{17_Gan2019_Seer}
        & \textsf{DT} 
        & \cellcolor[HTML]{97CDD8} supervised learning
        %& online
        & \textsf{P}
        & \textsf{S}
        & previous and runtime traces
        \\
    
    Nedelkoski \etal \cite{34_Nedelkoski2019_AnomalyDetectionDistributedTracing}
        & \textsf{DT}  %(with Zipkin
        & \cellcolor[HTML]{97CDD8} supervised learning
        %& online
        & \textsf{P}
        & \textsf{S}
        & previous and runtime traces
        \\
    
    Bogatinovski \etal \cite{36_Bogatinovski2020_SelfSupervisedAnomalyDetection}
        & \textsf{DT} 
        & \cellcolor[HTML]{97CDD8} supervised learning
        & \textsf{P}
        & \textsf{A}
        & previous and runtime traces
        \\
    
    MEPFL \cite{02_Zhou2019_LatentErrorPrediction}
        & \textsf{DT} 
        & \cellcolor[HTML]{97CDD8} supervised learning 
        %& online               
        & \textsf{F}
        & \textsf{S}
        & app deployment, test cases
        \\
    Meng \etal \cite{23_Meng2021_DetectingAnomaliesMicroservices}
        & \textsf{DT}  %(with Jaeger, Zipkin
        & \cellcolor[HTML]{FFD1C2} trace comparison
        %& offline % in TTV, authors state that the proposed approach is too costly to enact online detection on medium/large-scale multi-service apps
        & \textsf{F}, \textsf{P}
        & \textsf{S}
        & previous and runtime traces 
        \\
    Wang \etal \cite{26_Wang2020_WorkflowAwareFaultDiagnosis}
        & \textsf{DT}  %(with OpenTracing
        & \cellcolor[HTML]{FFD1C2} trace comparison
        %& offline % in TTV, authors state that the proposed approach is too costly to enact online detection on medium/large-scale multi-service apps
        & \textsf{F}, \textsf{P}
        & \textsf{S}
        & previous and runtime traces
        \\
    Chen \etal \cite{39_Chen2020_MatrixSketchBasedAnomalyDetection}
        & \textsf{DT} 
        & \cellcolor[HTML]{FFD1C2} trace comparison
        %& online
        & \textsf{P}
        & \textsf{S}
        & previous and runtime traces
        \\
        \thinline
    Gulenko \etal \cite{37_Gulenko2018_DetectingAnomalousBehaviourBlackBoxServices}
        & \textsf{M} 
        & \cellcolor[HTML]{E7F2F8} unsupervised learning
        %& online
        & \textsf{P}
        & \textsf{S}
        & app deployment, workload generator
        \\
    LOUD \cite{13_Mariani2018_LocalizingFaultsCloudSystems}
        & \textsf{M}  
        & \cellcolor[HTML]{E7F2F8} unsupervised learning
        %& online
        & \textsf{P}
        & \textsf{S}
        & app deployment, workload generator
        \\
    MicroRCA \cite{04_Wu2020_MicroRCA} 
        & \textsf{M} 
        & \cellcolor[HTML]{E7F2F8} unsupervised learning
        %& online
        & \textsf{P}
        & \textsf{S}
        & k8s app deployment, workload generator 
        \\ 
    
    Wu \etal \cite{09_Wu2020_MicroserviceDiagnosisDeepLearning}
        & \textsf{M} 
        & \cellcolor[HTML]{E7F2F8} unsupervised learning
        %& online
        & \textsf{P}
        & \textsf{S}
        & k8s app deployment, workload generator
        \\
    
    DLA \cite{14_Samir2019_DLA}
        & \textsf{M} 
        & \cellcolor[HTML]{E7F2F8} unsupervised learning
        %& online
        & \textsf{P}
        & \textsf{S}
        & k8s app deployment % HHMM model - not for anomaly detection!
        \\
    
    CloudRanger \cite{29_Wang2018_CloudRanger}
        & \textsf{M} 
        & \cellcolor[HTML]{E7F2F8} unsupervised learning
        %& online
        & \textsf{P}
        & \textsf{A}
        & app deployment, prev. monitored KPIs
        \\
    
    Hora \cite{27_Pitakrat2018_HoraExtended,28_Pitakrat2016_Hora}
        & \textsf{M} 
        & \cellcolor[HTML]{E7F2F8} unsupervised learning
        %& online
        & \textsf{P}
        & \textsf{S}
        & app deployment
        \\
    ADS \cite{15_Du2018_AnomalyDetectionContainerMicroservices}
        & \textsf{M} 
        & \cellcolor[HTML]{97CDD8} supervised learning
        & \textsf{P}
        & \textsf{S}
        & k8s app deployment, workload generator, failure injectors
        \\
    PreMiSE \cite{33_Mariani2020_PredictingFailures}
        & \textsf{M} 
        & \cellcolor[HTML]{97CDD8} supervised learning
        %& online
        & \textsf{F}, \textsf{P}
        & \textsf{S}
        & app deployment, workload generator  % with one service per VM
        \\
    
    CauseInfer \cite{40_Chen2020_CauseInfer,48_Chen2014_CauseInfer}
        & \textsf{M} 
        & \cellcolor[HTML]{F3EDCE} SLO check
        %& online
        & \textsf{P}
        & \textsf{A}
        & app deployment, SLOs
        \\
    
    MicroScope \cite{06_Lin2018_Microscope,19_Guan2019_MicroscopeDemo}
        & \textsf{M} %\textsf{M} 
        & \cellcolor[HTML]{F3EDCE} SLO check
        %& online
        & \textsf{P}
        & \textsf{A}
        & k8s app deployment, SLOs
        \\
    
    $\epsilon$-diagnosis \cite{49_Shan2019_epsilonDiagnosis}
        & \textsf{M} 
        & \cellcolor[HTML]{F3EDCE} SLO check
        %& online
        & \textsf{P}
        & \textsf{A}
        & k8s app deployment, SLOs
        \\
    M-MFSA- HDA \cite{24_Zang2018_FaultDiagnosisMicroservices}
        & \textsf{M} 
        & \cellcolor[HTML]{E6D6C7} heartbeating
        %& online
        & \textsf{F}
        & \textsf{S}
        & app deployment
        \\
    \thickline
\end{tabular}
\end{tiny}
\end{table}

\Cref{tab:detection} recaps the surveyed anomaly detection techniques by distinguishing their classes, \viz whether they are log-based, distributed tracing-based, or monitoring-based, as well as the method they apply to enact anomaly detection.
The table also classifies the surveyed techniques based on whether they can detect functional or performance anomalies, and on the \enquote{granularity} of detected anomalies, \viz whether they just state that an application is suffering from an anomaly, or whether they distinguish which of its services are suffering from anomalies.
Finally, the table provides insights on the artifacts that must be provided to the surveyed anomaly detection techniques, as input needed to actually enact anomaly detection.

Taking \Cref{tab:detection} as a reference, we hereafter summarise the surveyed anomaly detection techniques (\Cref{sec:detection:discussion:summary}).
We also discuss them under three different perspectives, \viz finding a suitable compromise between type/granularity of detected anomalies and setup costs (\Cref{sec:detection:discussion:type-granularity-costs}), their accuracy, especially in the case of dynamically changing applications (\Cref{sec:detection:discussion:adaptability}), and the need for explainability and countermeasures for detected anomalies (\Cref{sec:detection:discussion:explainability-countermeasures}). %\footnote{We focus on qualitatively comparing the surveyed anomaly detection techniques. A quantitative comparison (\eg of their accuracy and performance) is outside of the scope of this paper and left for future work (\Cref{sec:conclusions}).}

\subsubsection{Summary}
\label{sec:detection:discussion:summary}
All the surveyed techniques (but for those based on SLO checks \cite{40_Chen2020_CauseInfer,48_Chen2014_CauseInfer,19_Guan2019_MicroscopeDemo,06_Lin2018_Microscope,49_Shan2019_epsilonDiagnosis} or on heartbeating \cite{24_Zang2018_FaultDiagnosisMicroservices}) rely on processing data collected in training runs of the target applications. %, \viz the events they log, the traces they produce, or the KPIs monitored with monitoring agents.
The idea is to use the logs, traces, or KPIs collected in training runs of the application as a baseline against which to compare newly produced logs or traces, or newly monitored KPIs.
Machine learning is the most used approach to train baseline models of the application behaviour: the logs, traces, or KPIs collected during training runs of the application are indeed used to train baseline models with unsupervised learning algorithms
\cite{37_Gulenko2018_DetectingAnomalousBehaviourBlackBoxServices,31_Jia2017_AnomalyDiagnosisHybridGraphModel,32_Jia2017_LogSed,22_Jin2020_AnomalyDetectionMicroservicesRPCA,16_Liu2020_TraceAnomaly,13_Mariani2018_LocalizingFaultsCloudSystems,38_Nandi2016_AnomalyDetectionControlFlowGraphMining,35_Nedelkoski2019_AnomalyDetectionTracingData,28_Pitakrat2016_Hora,27_Pitakrat2018_HoraExtended,14_Samir2019_DLA,29_Wang2018_CloudRanger,04_Wu2020_MicroRCA,09_Wu2020_MicroserviceDiagnosisDeepLearning}, or with supervised learning algorithms if such information is also labelled with the failures that have happened or were injected on services during the training runs \cite{36_Bogatinovski2020_SelfSupervisedAnomalyDetection,15_Du2018_AnomalyDetectionContainerMicroservices,02_Zhou2019_LatentErrorPrediction,34_Nedelkoski2019_AnomalyDetectionDistributedTracing,33_Mariani2020_PredictingFailures,17_Gan2019_Seer}.
The trained baseline models are typically clusters defining the space where newly monitored KPIs should pertain to not be considered anomalous, or classifiers  or deep neural networks to be fed with newly monitored KPIs, logs, or traces to detect whether they denote anomalies.

An alternative approach to machine learning is trace comparison, which is enacted by Chen \etal \cite{39_Chen2020_MatrixSketchBasedAnomalyDetection}, Meng \etal \cite{23_Meng2021_DetectingAnomaliesMicroservices}, and Wang \etal \cite{26_Wang2020_WorkflowAwareFaultDiagnosis}, all based on instrumenting applications to feature distributed tracing.
The idea here is to store the traces that can possibly occur in an application, and then to check whether newly generated traces are similar enough to the stored ones.
This approach however results to be quite time consuming, indeed bringing two out of the three techniques enacting trace comparison to be suited only for offline anomaly detection.
Only Chen \etal \cite{39_Chen2020_MatrixSketchBasedAnomalyDetection} achieves time performances good enough to enact online anomaly detection on large-scale applications, whereas Meng \etal \cite{23_Meng2021_DetectingAnomaliesMicroservices} and Wang \etal \cite{26_Wang2020_WorkflowAwareFaultDiagnosis} explicitly state that their techniques are too time consuming to be used online on medium-/large-scale applications.

\subsubsection{Type and Granularity of Detected Anomalies vs. Setup Costs}
\label{sec:detection:discussion:type-granularity-costs}
The surveyed techniques achieve a different granularity in anomaly detection by applying different methods and requiring different artifacts (\Cref{tab:detection}).
Independently from the required instrumentation and from whether they apply machine learning or trace comparison, the techniques comparing newly collected KPIs, logs, or traces against a baseline turn out to be much finer in the granularity of detected anomalies, if compared with techniques enacting SLO checks.
Indeed, whilst monitoring SLOs on the frontend of an application allows to state that the whole application is suffering from a performance anomaly, the techniques based on machine learning or trace comparison often allow to detect functional or performance anomalies at service-level.
They indeed typically indicate which of its services are suffering from functional or performance anomalies, whilst also providing information on which are the KPIs monitored on a service, the events it logged, or the portion of a trace due to which a service is considered anomalous.
In the case of supervised learning-based techniques \cite{36_Bogatinovski2020_SelfSupervisedAnomalyDetection,15_Du2018_AnomalyDetectionContainerMicroservices,02_Zhou2019_LatentErrorPrediction,34_Nedelkoski2019_AnomalyDetectionDistributedTracing,33_Mariani2020_PredictingFailures,17_Gan2019_Seer}, they also indicate which failures may possibly correspond to the symptoms denoted by the detected anomalies.
All such information is precious for application operators as it eases further investigating on detected anomalies, \eg to check whether they correspond to services truly failing and to determine why they have failed \cite{Brogi2020_FailureCausalities}.

This comes at the price of requiring to execute training runs of multi-service applications to collect logs, traces, or KPIs to train the baseline models against which to compare those produced at runtime. 
In the case of supervised learning-based techniques, it is also required to either provide what needed to automatically inject failures or to label the collected data to distinguish whether it corresponds to normal runs or to runs where specific failures occurred on specific services.
% Log-based techniques directly work with the event logs produced by the application services, requiring those produced by all services in training runs of the application to train baseline models that are then compared against newly produced logs \cite{38_Nandi2016_AnomalyDetectionControlFlowGraphMining,31_Jia2017_AnomalyDiagnosisHybridGraphModel,32_Jia2017_LogSed}.
%They hence not require any application instrumentation, other than assuming that the application services log events, which is typically the case.
Then, log-based techniques does not require any application instrumentation, as they directly work on the event logs produced by the application services during the training phase and at runtime \cite{38_Nandi2016_AnomalyDetectionControlFlowGraphMining,31_Jia2017_AnomalyDiagnosisHybridGraphModel,32_Jia2017_LogSed}.
%
Distributed tracing-based techniques instead assume applications to be instrumented to feature distributed tracing \cite{Richardson2018_MicroservicesPatterns}, so as to automatically collect the traces of events produced by their services.
In this case, the traces used to train the baseline models are obtained twofold, \viz either only with failure-free runs of an application  \cite{16_Liu2020_TraceAnomaly,35_Nedelkoski2019_AnomalyDetectionTracingData,22_Jin2020_AnomalyDetectionMicroservicesRPCA,36_Bogatinovski2020_SelfSupervisedAnomalyDetection} or by also considering runs where specific failures occurred on specific services \cite{17_Gan2019_Seer,34_Nedelkoski2019_AnomalyDetectionDistributedTracing,02_Zhou2019_LatentErrorPrediction}.
%In the latter case, the possibilities to obtain both failure-free and failing traces are twofold, \viz using historically collected traces for an application by explicitly labelling those where failures are known to have occurred \cite{17_Gan2019_Seer}, or running the application by artificially injecting specific failures in specific services \cite{34_Nedelkoski2019_AnomalyDetectionDistributedTracing,02_Zhou2019_LatentErrorPrediction}.

Monitoring-based techniques work with another different setup: they require deploying monitoring agents together with the application services, and they directly work with the application services themselves (rather than processing their logs or the traces they produce).
Whilst this is the only need for SLO check-based techiniques \cite{40_Chen2020_CauseInfer,48_Chen2014_CauseInfer,49_Shan2019_epsilonDiagnosis} and for DLA \cite{14_Samir2019_DLA}, monitoring-based techniques typically require additional artifacts to generate a baseline modelling of the KPIs monitored on the application.
They indeed either require KPIs already monitored in former production runs of the application \cite{29_Wang2018_CloudRanger}, or workload generators to simulate the traffic that will then be generated by end users \cite{37_Gulenko2018_DetectingAnomalousBehaviourBlackBoxServices,04_Wu2020_MicroRCA,09_Wu2020_MicroserviceDiagnosisDeepLearning,33_Mariani2020_PredictingFailures,15_Du2018_AnomalyDetectionContainerMicroservices}.
ADS \cite{15_Du2018_AnomalyDetectionContainerMicroservices} also requires application operators to provide failure injection modules to inject specific failures in specific services during the training of the baseline model. 
In addition, some of the monitoring-based techniques focus on specific setups for the deployment of an application, \eg requiring their k8s deployment \cite{49_Shan2019_epsilonDiagnosis,04_Wu2020_MicroRCA,09_Wu2020_MicroserviceDiagnosisDeepLearning,14_Samir2019_DLA,06_Lin2018_Microscope,19_Guan2019_MicroscopeDemo} or assuming their VM-based deployment to be such that each service is deployed in a separate VM \cite{33_Mariani2020_PredictingFailures}.

% Among monitoring-based techniques, a quite much different setup is instead required by M-MFSA-HDA \cite{24_Zang2018_FaultDiagnosisMicroservices} to enact self-adaptive heartbeating to detect functionally anomalous services.
% M-MFSA-HDA \cite{24_Zang2018_FaultDiagnosisMicroservices} indeed not only requires deploying a monitoring agent, but also to instrument applications to feature an \enquote{ad-hoc} heartbeating protocol, \viz to allow services to promptly answer to the heartbeat messages sent by the monitoring agent.
% Such an \enquote{ad-hoc} instrumentation is not required by the other monitoring techniques, nor by those based on distributed tracing.
% Indeed, the other monitoring-based techniques require to install agents to monitor KPIs on application services, either by deploying external agents or by adapting the deployment of an application to be a service mesh, \eg with Istio \cite{Istio} and Prometheus \cite{Prometheus}.
% Distributed tracing techniques instead require to instrument applications to feature distributed tracing, \eg with Zipkin \cite{Zipkin} or OpenTracing \cite{OpenTracing}.
% All such types of instrumentation (but for heartbeating) are often already featured in modern multi-service applications, especially in microservices \cite{Richardson2018_MicroservicesPatterns,Soldani2018_MicroservicesPainsGains}.

In summary, not all discussed techniques are suitable to be used on existing applications as they are.
They may require to instrument multi-service applications, to adapt their deployment to work with a given technique, or to provide additional artifacts to configure a technique to enact anomaly detection on an application.
Hence, whilst some techniques are known to detect functional/performance anomalies at a finer granularity, this may not be enough to justify the cost for setting them up to work with a given application, \viz to suitably instrument the application or its deployment, or to provide the needed artifacts.
In some other cases, an application may already natively provide what is needed to apply it a given anomaly detection techniques, \eg event logging or distributed tracing.
The choice of which technique to use for enacting anomaly detection on a given application hence consists of finding a suitable trade-off between the desired granularity/type of detected anomalies and the cost for setting up such technique to work with the given application. 
We believe that the classification provided in this paper can provide a first support for application operators needing to take such a choice.

%\subsubsection{Adaptability to Varying Running Conditions}
\subsubsection{Accuracy of Anomaly Detection}
\label{sec:detection:discussion:adaptability}
The surveyed techniques (but for those based on SLO checks or heartbeating) share a common assumption, \viz that, at runtime, the application and its services behave similarly to the training runs.
In other words, they assume that applications will feature similar performances, log similar events, or perform similar interactions in a similar order. 
This is however not always the case, since the actual behaviour of multi-service applications also depends on the conditions under which they are executed, which often dynamically changes \cite{Soldani2018_MicroservicesPainsGains}.
For instance, if hosted on the same virtual environments but with different sets of co-hosted applications, multi-service applications may suffer from different contentions on computing resources and hence feature different performances \cite{Medel2018_ResourceManagementPerformanceKubernetes}.
Similar considerations derive from the workload generated by end-users, which often varies with seasonalities \cite{33_Mariani2020_PredictingFailures}.
In addition, multi-service applications can be highly dynamic in their nature, especially if based on microservices: new services can be included to feature new functionalities, or existing services can be replaced or removed as becoming outdated \cite{Soldani2018_MicroservicesPainsGains}.
As a result, all techniques assuming that multi-service applications keep running under the same conditions as in training runs may lower their accuracy if the running conditions of an application change, as they may end up with detecting false positives or with suffering from false negatives.
They may indeed detect functional or performance anomalies in a running application, even if these are not anomalous changes, but rather correspond to the \enquote{new normal behaviour} in the new conditions under which an application runs.
They may instead consider the actual behaviour of an application as normal, even if such behaviour may denote anomalies in the latest running conditions of an application.

The problem of false positives and false negatives is not to be underestimated.
Anomaly detection is indeed enacted as anomalies are possible symptoms of failures.
A detected anomaly on a service hence provides a warning on the fact that such service can have possibly failed, and it may also be used to determine the type of failure affecting such service. 
A false positive, \viz a detected anomaly that is actually not an anomaly, would result the application operator in spending resources and efforts on application services that could have possibly failed, \eg to recover them or to understand why they have failed, even if this was not the case. 
Even worse is the case of false negatives: if some anomalies would not be detected, the symptoms of some failures may not get detected, which in turn means that the application operator would not be warned in case such failures actually occur in a multi-service application.

To mitigate this issue, the baseline used to detect anomalies should be kept up to date by adapting it to the varying conditions under which a multi-service application runs.
A subset of the surveyed techniques explicitly recognise this problem, \eg PreMiSE \cite{33_Mariani2020_PredictingFailures} prescribes to use a huge amount of training data to include as many evolutions of the application as possible, therein including seasonalities of workloads, to mitigate the effects of the \enquote{training problem} as much as possible.
Some of the surveyed techniques also propose solutions to address the training problem by keeping the baseline model up to date.
For instance, DLA \cite{14_Samir2019_DLA}, Seer \cite{17_Gan2019_Seer}, and Wang \etal \cite{26_Wang2020_WorkflowAwareFaultDiagnosis}, among others, propose to periodically re-train the baseline model by considering newly collected data as training data.
They also observe that the re-training period should be tailored as a trade-off between the effectiveness of the trained models and the computational effort needed to re-train the model, as re-training is typically resource and time consuming. 
CloudRanger \cite{29_Wang2018_CloudRanger} instead keeps the baseline continuously updated by enacting continuous learning on the services forming an application.
Among trace comparison-based techniques, Chen \etal \cite{39_Chen2020_MatrixSketchBasedAnomalyDetection} actually try to address the issue with false positives by exploiting matrix sketching to keep track of the entire prior for each service.
They hence compare newly collected information with all what has been observed formerly, both during the offline reference runs and while the application was running in production.

In the case of supervised learning-based techniques \cite{36_Bogatinovski2020_SelfSupervisedAnomalyDetection,15_Du2018_AnomalyDetectionContainerMicroservices,02_Zhou2019_LatentErrorPrediction,34_Nedelkoski2019_AnomalyDetectionDistributedTracing,33_Mariani2020_PredictingFailures,17_Gan2019_Seer}, the false positives and false negatives possibly caused by the training problem impact not only on anomaly detection, but also on the identification of the failure corresponding to a detected anomaly. 
Failure identification is indeed based on the correlation of symptoms, assuming that the same failure happening on the same service would result in a similar set of symptoms. 
Again, the varying conditions under which an application runs, as well as the fact that the service forming an application may dynamically change themselves, may result in the symptoms observed on a service to change over time.
The failures that can possibly affect a service may also change over time, hence requiring to adapt the supervised learning-based techniques not only to the varying conditions under which an application runs, but also to the varying sets of possible failures for a service.
To tackle this issue, Seer \cite{17_Gan2019_Seer} re-trains the baseline modelling of the application by also considering the failures that were observed on the application services while they were running in production.
Seer \cite{17_Gan2019_Seer} is actually the only technique  dealing with the training problem among those enacting supervised learning.
This is mainly because Seer \cite{17_Gan2019_Seer} is designed to be periodically trained on the data monitored while an application is running in production, whereas the other supervised learning-based techniques rely on predefined failures being artificially injected during the training runs.

The above however only constitute mitigation actions, \viz actions allowing to mitigate the false positives/negatives due to the training problem. 
False positives/negatives are indeed a price to pay when detecting anomalies by comparing the runtime behaviour of an application with a baseline modelling of its normal, failure-free behaviour.
A quantitative comparison of the accuracy of the different techniques on given applications in given contexts would further help application operators in this direction, and it will complement the information in this survey in supporting application in choosing the anomaly detection techniques best suited to their needs.
Such a quantitative comparison is however outside of the scope of this survey and a direction for future work.

\subsubsection{Explainability and Countermeasures}
\label{sec:detection:discussion:explainability-countermeasures}
The problem of false positives/negatives is also motivated by a currently open challenge in the field of anomaly detection for multi-service applications, \viz \enquote{explainability}.
This is tightly coupled with the problem of explainable AI \cite{Guidotti2019_Explainability}, given that anomaly detection is mainly enacted with machine learning-based techniques.
If giving explanations for answers given by learned models is generally important, it is even more important in online anomaly detection.
Explanations would indeed allow application operators to directly exclude false positives, or to exploit such explanations, \eg to investigate the root causes for the corresponding failures, as well as to design countermeasures avoiding such failures to happen again \cite{Brogi2020_FailureCausalities}.
Root cause analysis techniques do exist and ---as we show in \Cref{sec:rca}--- they still suffer from the training problem potentially causing false positives or false negatives, as well as from the need for explainability.

Another interesting research direction would be to take advantage of the information processed to detect anomalies to enact countermeasures to isolate (the failures causing) detected anomalies,
\eg by proactively activating circuit breakers/bulkheads \cite{Richardson2018_MicroservicesPatterns} to avoid that a service anomaly propagates to other services.
%
%Whilst avoiding (the failure corresponding to) a detected anomaly of course needs acting on the root causes for such an anomaly, countermeasures could anyhow be taken to prevent the anomaly in a service from propagating to other services, \eg by activating circuit breakers/bulkheads \cite{Richardson2018_MicroservicesPatterns} to avoid failure propagation.
%None of the surveyed techniques however exploits detected anomalies, nor the corresponding logs, traces, or monitored KPIs, to suggest countermeasures to the application operator.
A first attempt in this direction is made 
by Seer \cite{17_Gan2019_Seer}, which notifies the system manager about the anomaly detected on a service and the failure that can have possibly caused it, \eg CPU hog or memory leak. 
With such information, Seer \cite{17_Gan2019_Seer} envisions the system manager to be able to enact policies mitigating the effects of the failure causing a detected anomaly, \eg by providing the node hosting an anomalous service with additional computing resources.